% =====================================================================
%  Methodology chapter -- AI-Assisted Rooftop Solar Potential Estimation
%  and Financial Analysis System.
%
%  This file is a FRAGMENT, intended to be \input{} from a master
%  thesis file (e.g. bachelor.tex) that provides the document class,
%  preamble, \begin{document}, \bibliography{references}, etc.
%
%  Required packages in the master preamble:
%      \usepackage[utf8]{inputenc}     % only with pdflatex
%      \usepackage[T1]{fontenc}
%      \usepackage{amsmath, amssymb, mathtools}
%      \usepackage{booktabs}
%      \usepackage{tabularx}
%      \usepackage{longtable}
%      \usepackage{enumitem}
%      \usepackage{listings}
%      \usepackage{hyperref}
%      \usepackage[round,authoryear]{natbib}   % or {numbers,sort&compress}
% =====================================================================

% -- Local macros (safe if the master file already defines them) -------
\providecommand{\degC}{\,^{\circ}\mathrm{C}}
\providecommand{\COtwo}{\mathrm{CO_{2}}}

% =====================================================================
\chapter{Methodology}
\label{ch:methodology}
% Replace \chapter with \section if your thesis class is `article`.
% =====================================================================

\noindent\textbf{Status:} Day-19 deliverable --- academic, LaTeX-ready.\\
\textbf{Pairs with:} \texttt{validation.tex} (Day 20),
\texttt{limitations.tex} (Day 21), \texttt{references.bib} (Day 21).
All citations below resolve against \texttt{references.bib}.

\bigskip

This chapter documents the analytical and computational methodology
behind the \emph{AI-Assisted Rooftop Solar Potential Estimation and
Financial Analysis System}. Each section maps one-to-one to a backend
service module under \texttt{backend/app/services/}, so an examiner
can read the chapter top-to-bottom and audit each step against the
Python source. Equations are written in \LaTeX{} so the chapter is
drop-in ready for the thesis manuscript without re-typing.

% =====================================================================
\section{Data Sources}
% =====================================================================

\subsection{Solar irradiance --- PVGIS Typical Meteorological Year}

The system consumes hourly typical-meteorological-year (TMY) records
from the European Commission's Photovoltaic Geographical Information
System (PVGIS) v5.2 \citep{pvgis}, retrieved via the public REST
endpoint \url{https://re.jrc.ec.europa.eu/api/v5_2/}. For every site
$(\varphi, \lambda)$ supplied as latitude / longitude in decimal
degrees, PVGIS returns 8\,760 hourly tuples
\[
\bigl\{(\mathrm{GHI}_h, \mathrm{DNI}_h, \mathrm{DHI}_h, T_{\text{air},h},
v_{\text{wind},h})\bigr\}_{h=1}^{8760}
\]
in $\mathrm{W\cdot m^{-2}}$, ${}^{\circ}\mathrm{C}$, and
$\mathrm{m\cdot s^{-1}}$ respectively. The TMY series is constructed
by PVGIS's typical-month selection algorithm, applied to ERA-5
reanalysis aggregated over the 2005--2020 period at $\sim$5\,km grid
resolution.

A TMY is, by construction, a \emph{typical} climatic year --- not a
forecast of any specific calendar year. A single PV simulation
against a TMY therefore captures the \emph{expected} generation under
climatologically representative weather. Year-to-year weather
variability is restored explicitly in \S\ref{sec:montecarlo} via the
Monte Carlo \texttt{annual\_yield\_factor} distribution.

\subsection{Roof geometry --- OpenStreetMap + Google Maps Static}

Building footprints come from the OpenStreetMap Overpass API
\citep{osm-overpass}. For a query point $(\varphi_0, \lambda_0)$ the
client retrieves every closed \texttt{way} tagged \texttt{building=*}
whose centroid lies within the configured search radius
$r_{\text{search}} = 50\,\mathrm{m}$ (Egypt-tuned default; bounded
above at $500\,\mathrm{m}$ to protect the public Overpass instance).
Satellite raster tiles for the same coordinate come from the Google
Maps Static API at zoom 20 with \texttt{scale=2}, giving an effective
ground resolution of ${\sim}0.14~\mathrm{m\cdot pixel^{-1}}$ at
Cairo's latitude --- sufficient for roof-edge gradient extraction
(\S\ref{sec:cv}).

\subsection{Egypt-specific operating constants}
\label{sec:constants}

\begin{table}[!htbp]
\centering
\footnotesize
\setlength{\tabcolsep}{4pt}
\renewcommand{\arraystretch}{1.15}
\begin{tabularx}{\linewidth}{@{}l X r X@{}}
\toprule
\textbf{Symbol} & \textbf{Description} & \textbf{Value} & \textbf{Source} \\
\midrule
$\eta_{\text{inv}}$         & Inverter efficiency (constant)        & 0.96
  & Modern grid-tied c-Si default; IRENA \citep{irena2023} \\
$f_{\text{util}}$           & Roof utilisation factor               & 0.70
  & Egyptian rooftop PV pre-feasibility \citep{mahmoud2023} \\
$\beta_{\text{def}}$        & Default tilt                          & $26^{\circ}$
  & Cairo latitude optimum \\
$A_{\text{def}}$            & Default azimuth                       & $180^{\circ}$
  & Northern-hemisphere optimum \\
$\eta_{\text{soil}}$        & Soiling loss range                    & 2--8\,\%
  & Egyptian PV literature \\
$\varepsilon_{\text{grid}}$ & Grid emission factor
  & 0.46
  & EEHC 2022/2023 \citep{eehc_emissions}, $\mathrm{kg\,\COtwo/kWh}$ \\
$C_{\text{kW}}$             & Installed cost                        & 35\,000
  & Egyptian solar market 2024, EGP/kW \\
$r_{\text{disc}}$           & Real discount rate                    & 0.04
  & Standard project finance \\
$i_{\text{tar}}$            & Tariff escalation (mean)              & 0.08
  & EgyptERA decade trend \citep{egyptera_tariff}, $\mathrm{yr^{-1}}$ \\
$d$                         & Module degradation rate               & 0.005
  & NREL median for mono-Si \citep{jordan2013}, $\mathrm{yr^{-1}}$ \\
$f_{\text{O\&M}}$           & Annual O\&M as fraction of capex      & 0.01
  & IRENA residential rooftop \citep{irena2023} \\
\bottomrule
\end{tabularx}
\caption{Egypt-specific operating constants used by the deterministic kernels.}
\label{tab:constants}
\end{table}

Every constant in this table is exposed as a Pydantic-validated
override on the corresponding endpoint, so a methodology-aware caller
can substitute a site-specific value without touching the source.

% =====================================================================
\section{PV Modelling}
\label{sec:pv}
% =====================================================================

\subsection{System sizing}

The sizing kernel converts a roof area $A_{\text{roof}}$ in
$\mathrm{m^{2}}$ into a discrete panel count $N$ and a nameplate DC
capacity $P_{\text{sys}}$ in kW:
\begin{equation}
A_{\text{usable}} = f_{\text{util}} \cdot A_{\text{roof}}, \qquad
N = \left\lfloor \frac{A_{\text{usable}}}{A_{\text{panel}}} \right\rfloor,
\qquad
P_{\text{sys}} = \frac{N \cdot P_{\text{panel}}}{1000}~[\mathrm{kW}].
\label{eq:sizing}
\end{equation}

Three design choices warrant explicit justification:

\begin{enumerate}[leftmargin=*]
\item \textbf{Floor, not round.} A fractional panel is physically
meaningless; rounding \emph{down} is the conservative,
installer-realistic choice that prevents an upward bias from
propagating through every downstream financial metric.
\item \textbf{Lumped utilisation factor.} The factor
$f_{\text{util}}=0.70$ absorbs edge setbacks (Egyptian fire and
maintenance codes), inter-row walkways for cleaning (a relevant loss
in Cairo's high-soiling environment), inter-row shading spacing for
tilted mounts at $\beta=26^{\circ}$, and roof-mounted obstructions
(water tanks, satellite dishes, parapets, HVAC condensers) without
explicit geometric modelling. The factor is the central value
reported by Egyptian rooftop PV pre-feasibility studies
\citep{mahmoud2023}.
\item \textbf{Hardware defaults reflect the 2024 Egyptian residential
market.} $P_{\text{panel}} = 450~\mathrm{W}$ on
$A_{\text{panel}} = 1.8~\mathrm{m^{2}}$ is the dominant module
configuration; the kernel echoes the assumptions in the response so
the JSON itself is self-auditing.
\end{enumerate}

\subsection{Energy generation --- pvlib model (industry-standard reference)}
\label{sec:pvlib}

The reference energy chain is implemented with \texttt{pvlib} Python
\citep{pvlib}, using the individual building blocks rather than the
high-level \texttt{ModelChain} so each step is explicit and traceable
in this chapter:

\begin{enumerate}[leftmargin=*]
\item \textbf{Solar geometry.}
$\theta_z, \gamma_s = \mathrm{NREL\,SPA}(\varphi, \lambda, t)$
via \texttt{pvlib.solarposition.get\_solarposition}.

\item \textbf{Plane-of-array (POA) irradiance.} Hay--Davies
sky-diffuse transposition (\texttt{pvlib.irradiance.get\_total\_irradiance}):
\begin{equation}
G_{\text{POA}} = G_{\text{beam}} \cos\theta_i
+ G_{\text{diff}} R_d^{\text{HD}}(\theta_z, \gamma_s, \beta)
+ \rho \, G_{\text{horiz}} \frac{1-\cos\beta}{2},
\label{eq:poa}
\end{equation}
with surface tilt $\beta$, surface azimuth $\gamma_p$, ground albedo
$\rho=0.20$ (PVWatts default for mixed urban surfaces).

\item \textbf{Cell temperature.} SAPM open-rack glass-polymer
parameter set (\texttt{pvlib.temperature.sapm\_cell}). The thermal
model is the standard choice for free-standing rooftop installations
of crystalline-silicon modules --- the dominant technology in the
Egyptian market.

\item \textbf{DC power.} Linear PVWatts equation
(\texttt{pvlib.pvsystem.pvwatts\_dc}):
\begin{equation}
P_{\text{dc}} = \frac{G_{\text{POA}}}{1000~\mathrm{W/m^{2}}}
\, P_{\text{dc0}}
\, \bigl[1 + \gamma_{P}(T_{\text{cell}} - 25\,^{\circ}\mathrm{C})\bigr],
\label{eq:pvwatts}
\end{equation}
with $\gamma_{P} = -0.0035\,/\,^{\circ}\mathrm{C}$ (typical for
monocrystalline silicon) \citep{pvwatts2014}.

\item \textbf{System DC losses.} A single combined-losses factor
folds in soiling, mismatch, wiring, and module-nameplate tolerance.
PVWatts canonical default $\eta_{\text{loss}} = 0.14$:
\begin{equation}
P_{\text{dc, net}} = (1 - \eta_{\text{loss}}) \, P_{\text{dc}}.
\label{eq:dcloss}
\end{equation}

\item \textbf{AC conversion.} \texttt{pvlib.inverter.pvwatts} clips
at the inverter nominal AC capacity using the constant efficiency
$\eta_{\text{inv}}=0.96$; the DC\,:\,AC ratio is held at unity in the
baseline.
\end{enumerate}

The annual energy, monthly energy, capacity factor, performance
ratio, and specific yield are aggregated from the resulting hourly
$P_{\text{ac}}$ series in the standard textbook fashion. The
``Cairo sanity band'' for the headline specific yield is
$1\,700 \le P_{\text{ac, ann}} / P_{\text{sys}} \le 1\,900~\mathrm{kWh/kWp}$;
agreement within this band is the first internal validation gate
(Validation \S1, Day 20).

\subsection{Energy generation --- Manual physics model (independent twin)}
\label{sec:manual}

The methodology's first defensible academic contribution is the
\textbf{dual-energy backbone}: alongside \texttt{pvlib} we ship a
from-first-principles chain that uses no \texttt{pvlib} calls. The
two chains share inputs and outputs but disagree, by design, at two
intermediate steps; cross-validating their headline numbers tells the
reader how much of the predicted yield is robust across modelling
assumptions and how much is parameter-sensitive.

\begin{table}[!htbp]
\centering
\footnotesize
\setlength{\tabcolsep}{4pt}
\renewcommand{\arraystretch}{1.2}
\begin{tabularx}{\linewidth}{@{}l X X@{}}
\toprule
\textbf{Step} & \textbf{pvlib chain} & \textbf{Manual chain} \\
\midrule
Solar geometry    & NREL SPA \citep{nrel-spa}
                  & Cooper (1969) declination + Spencer (1971) equation of time + classical hour-angle formulas \citep{duffie} \\
Sky-diffuse model & Hay--Davies (anisotropic)
                  & Liu--Jordan (isotropic) \citep{liu1960} \\
Cell temperature  & SAPM open-rack glass-polymer
                  & NOCT model:
                    $T_{\text{cell}} = T_{\text{air}} + \frac{NOCT-20}{800} \cdot G_{\text{POA}}$,
                    $NOCT = 45^{\circ}\mathrm{C}$ \\
DC power          & Linear PVWatts
                  & Same equation, parameter-identical $\gamma_{P}$ \\
DC losses         & PVWatts canonical $\eta_{\text{loss}} = 0.14$
                  & Same \\
AC conversion     & \texttt{pvlib.inverter.pvwatts}
                  & Constant $\eta_{\text{inv}}$, clip at $P_{\text{sys}}$ \\
\bottomrule
\end{tabularx}
\caption{Deliberate methodological disagreements between the pvlib
and manual energy chains.}
\label{tab:dual}
\end{table}

The two intermediate disagreements (Hay--Davies vs Liu--Jordan, SAPM
vs NOCT) are deliberate. They span the methodological choices a
practitioner would actually make; the residual between the two
annual energies is therefore an honest \emph{methodological-uncertainty}
number rather than an artefact of two implementations of the same
paper.

The cross-validation classification used in the dashboard is
\begin{equation}
\Delta_{\text{rel}} =
  \frac{E_{\text{manual}} - E_{\text{pvlib}}}{E_{\text{pvlib}}},
\qquad
\text{verdict} =
\begin{cases}
\text{strong agreement}     & |\Delta_{\text{rel}}| < 5\,\% \\
\text{reasonable agreement} & 5\,\% \le |\Delta_{\text{rel}}| < 10\,\% \\
\text{material divergence}  & |\Delta_{\text{rel}}| \ge 10\,\%
\end{cases}
\label{eq:verdict}
\end{equation}

The 5\,\% strong-agreement threshold matches the documented
pvlib-vs-PVSyst band for c-Si rooftop installations \citep{pvlib}.

\subsection{Loss decomposition}

The DC-side loss factor $\eta_{\text{loss}} = 0.14$ used by both
energy chains decomposes into the components published in the NREL
PVWatts Version 5 manual \citep{pvwatts2014}:

\begin{table}[!htbp]
\centering
\footnotesize
\begin{tabular}{@{}lr@{}}
\toprule
\textbf{Component} & \textbf{Loss} \\
\midrule
Soiling                          & 2\,\%   \\
Shading                          & 3\,\%   \\
Snow                             & 0\,\% (Egypt) \\
Mismatch                         & 2\,\%   \\
Wiring                           & 2\,\%   \\
Connections                      & 0.5\,\% \\
Light-induced degradation (LID)  & 1.5\,\% \\
Nameplate rating                 & 1\,\%   \\
Availability                     & 3\,\%   \\
\midrule
\textbf{Total (combined)}        & $\mathbf{\approx 14\,\%}$ \\
\bottomrule
\end{tabular}
\caption{NREL PVWatts canonical DC-side loss decomposition.}
\label{tab:losses}
\end{table}

The Egypt context would usually argue for a higher soiling allowance
than 2\,\% --- Cairo's $\mathrm{PM}_{10}$ load is ${\sim}3\times$ the
European average, and field studies have reported 5--8\,\% soiling
losses during the long dry season \citep{elsayed2017}. The 14\,\%
default is therefore a deliberately conservative \emph{floor}; the
Monte Carlo engine (\S\ref{sec:montecarlo}) treats the lumped loss as
part of the \texttt{annual\_yield\_factor} distribution, so both the
headline number and the uncertainty interval reflect the realistic
spread.

% =====================================================================
\section{Roof Detection (Contribution A)}
\label{sec:roof}
% =====================================================================

The first thesis contribution is to \textbf{eliminate the roof area
question}: a homeowner who does not know their roof area to within
$\pm 10~\mathrm{m^{2}}$ would otherwise propagate that uncertainty
straight through to the headline payback. The detection pipeline is
two-stage by design: a vector OSM stage (\S\ref{sec:osm}) gives an
analytically defensible footprint for every successfully tagged
building, and a raster CV stage (\S\ref{sec:cv}) refines the polygon
and assigns a confidence score. A satellite tile that cannot be
loaded triggers a graceful fall-through to OSM-only with
$\text{conf}=0$, so the pipeline always produces a result.

\subsection{OSM footprint retrieval and selection (Day 10)}
\label{sec:osm}

For a query point $(\varphi_0, \lambda_0)$ the orchestrator issues an
Overpass \texttt{way[building]} query within the configured search
radius and ranks the returned polygons by a deterministic two-stage
rule:

\begin{enumerate}[leftmargin=*]
\item \textbf{Containment first.} Even a 1\,cm offset from a building
edge is far stronger evidence of identity than centroid distance ---
a user who drops a pin on their roof is almost always inside the
polygon.
\item \textbf{Smallest area among containers.} Egyptian residential
parcels are sometimes nested inside a larger
\texttt{building=apartments} polygon that wraps a courtyard and
several units; picking the \emph{innermost} container disambiguates
correctly.
\end{enumerate}

When \textbf{no} polygon contains the pin (rural plot, mis-mapped
building, or sparse OSM coverage), the orchestrator falls back to the
building with the closest centroid; the result schema flags
\texttt{contains\_query\_point = False} so the dashboard can warn the
user that the match is heuristic.

For local distance and area we use a small-region equirectangular
projection centred on the user's pin:
\begin{equation}
x = (\lambda - \lambda_0) \cos\varphi_0 \cdot R_{\oplus}
       \cdot \frac{\pi}{180},
\qquad
y = (\varphi - \varphi_0) \cdot R_{\oplus} \cdot \frac{\pi}{180},
\label{eq:projection}
\end{equation}
with $R_{\oplus} = 6\,378\,137~\mathrm{m}$ (WGS84 semi-major axis).
At the building scale ($< 200~\mathrm{m}$ diagonal) the maximum area
distortion versus a UTM-quality projection is below 0.05\,\% at
Cairo's latitude --- well inside the roof-utilisation-factor
uncertainty downstream. We deliberately avoid pulling in
\texttt{pyproj} to keep the dependency surface small.

\subsection{CV regularisation, confidence, and orientation (Day 11)}
\label{sec:cv}

The OSM polygon is digitised by hand and routinely off by 1--2\,m at
corners \citep{osm-quality}. The Day-11 refinement combines the OSM
prior with the satellite tile to produce two outputs:

\paragraph{Polygon regularisation.}
Real residential rooftops are almost always rectangular, and OSM
contributors often capture noisy traces of an essentially rectangular
building. The kernel computes the \emph{minimum-rotated rectangle} of
the OSM polygon (\texttt{shapely}'s
\texttt{minimum\_rotated\_rectangle}) and emits that as the refined
geometry. A regularised polygon collapses the digitisation noise back
to the underlying rectangle so the energy model receives a clean
shape and a single dominant orientation.

\paragraph{Confidence score.}
A polygon's segmentation confidence is the \emph{fraction} of
polygon-perimeter pixels whose local Sobel gradient magnitude exceeds
the image's median, normalised to $[0, 1]$. Roof--ground transitions
show up as strong intensity gradients (concrete vs ground, shingle vs
lawn), so a polygon whose edges trace those gradients gets a high
score; one that floats inside a uniform region gets a low one.
Concretely, for an image $I$ with Sobel magnitude $|\nabla I|$ and
median $m_{\nabla}$, and a polygon perimeter sampled into pixel-band
coordinates $\mathcal{B}$ of width $w_b = 1.0~\mathrm{m}$
($\approx 7$\,px at zoom 20 / scale 2),
\begin{equation}
\text{conf} = \frac{1}{|\mathcal{B}|}
  \sum_{(u,v) \in \mathcal{B}}
  \mathbb{1}\bigl[|\nabla I|(u, v) > m_{\nabla}\bigr] \in [0, 1].
\label{eq:confidence}
\end{equation}
The 1.0\,m band width is wide enough to absorb GPS jitter and OSM
digitisation noise (typical 0.5--1.5\,m positional error per OSMF
accuracy studies \citep{osm-quality}) without bleeding into
neighbouring rooftops.

\paragraph{Orientation inference.}
The polygon's principal axis (long edge of the regularised rectangle)
gives a panel azimuth estimate modulo $180^{\circ}$. If the long-edge
azimuth is within $\delta_{\text{snap}} = 8^{\circ}$ of a cardinal
heading, the kernel \emph{snaps} to that heading --- installers
almost always orient panels along the building's primary edges, and
OSM digitisation jitter can rotate a square footprint by
1--2$^{\circ}$. The tilt is taken from OSM \texttt{roof:angle} tags
when present (validated against $0^{\circ} \le \beta \le 60^{\circ}$),
else inferred from \texttt{roof:shape}: flat for a flat slab (the
dominant Egyptian residential roof type, in which case the panel-tilt
prior reverts to the latitude optimum $\beta = 26^{\circ}$),
$30^{\circ}$ for pitched, $15^{\circ}$ for shed.

\paragraph{Why classical CV over a deep segmenter?}
A pre-trained Mask R-CNN or SAM model would push the headline number
on benchmark datasets, but
\textbf{(i)} it would add ${\sim}500~\mathrm{MB}$ of model weights,
GPU dependence, and a reproducibility burden incompatible with a
one-laptop thesis demo;
\textbf{(ii)} the marginal gain on dense urban Egyptian rooftops ---
already well-represented in OSM --- is bounded by the OSM positional
accuracy itself (0.5--1.5\,m), which classical regularisation already
matches;
\textbf{(iii)} confidence calibration of CNN segmenters in the
Egyptian residential domain has not been published and would be a
thesis in its own right \citep{vargas-munoz2021}.

% =====================================================================
\section{Tiered-tariff billing and PV-size optimisation (Contribution B)}
\label{sec:tariff}
% =====================================================================

\subsection{EgyptERA progressive marginal block tariff}
\label{sec:tariff-blocks}

Egypt's residential bill is a \emph{progressive marginal block}
tariff, not a flat rate. For the post-July 2023 EgyptERA reform
schedule \citep{egyptera_tariff}, the upper kWh/month bound and
EGP/kWh price for each tier $i = 1 \dots K$ are

\begin{table}[!htbp]
\centering
\footnotesize
\begin{tabular}{@{}crr@{}}
\toprule
$i$ & \textbf{Upper bound (kWh / mo)} & \textbf{Price (EGP / kWh)} \\
\midrule
1 & 50      & 0.58 \\
2 & 100     & 0.68 \\
3 & 200     & 0.83 \\
4 & 350     & 1.25 \\
5 & 650     & 1.40 \\
6 & 1\,000  & 1.45 \\
7 & $\infty$ & 1.55 \\
\bottomrule
\end{tabular}
\caption{Post-July 2023 EgyptERA residential tariff schedule.}
\label{tab:tariff}
\end{table}

For a household with monthly consumption $C~\mathrm{kWh/mo}$, with
$U_0 = 0$ and $U_i$ the upper bound of tier $i$, the
\emph{progressive marginal} monthly bill is
\begin{equation}
B(C) = \sum_{i=1}^{K} p_i \cdot
       \min\bigl(U_i - U_{i-1},\, \max(0,\, C - U_{i-1})\bigr).
\label{eq:bill}
\end{equation}

\paragraph{Why marginal, not inclusive?}
Egypt's published tariff is widely reported in two forms in
consumer-facing press: \emph{inclusive} (exceeding a threshold
reverts the \emph{whole} month's consumption to the higher tier) and
\emph{marginal} (each band charges only the kWh that fall inside it
--- the canonical ``progressive'' tariff used in most economics
texts). The official EgyptERA bill statement uses the marginal
interpretation (tier amounts on the bill add up to the total), and it
is the conservative choice for the \emph{PV value} claim: under an
inclusive schedule PV would look even \emph{more} valuable, since
saving one kWh near a band edge could shift the whole month down a
tier. Choosing the marginal schedule means our payback claims hold
under the worst of the two reasonable readings.

\subsection{Self-consumption-first PV netting}

The kernel models the \emph{self-consumption-first} metering
convention, which matches Egypt's residential default (no export
credit). For a household with monthly consumption $C_m$ and PV
generation $G_m$, the post-PV monthly consumption is
\begin{equation}
C^{\text{PV}}_m = \max\bigl(0,\, C_m - G_m\bigr),
\label{eq:netting}
\end{equation}
and the surplus generation $\max(0, G_m - C_m)$ is awarded an
optional export credit of $p_{\text{exp}}$ EGP/kWh (default 0).
Annual savings are
\begin{equation}
S_{\text{ann}} = \sum_{m=1}^{12}
  \bigl[ B(C_m) - B(C^{\text{PV}}_m)
       + p_{\text{exp}} \cdot \max(0,\, G_m - C_m) \bigr].
\label{eq:savings}
\end{equation}

The dashboard's headline \texttt{average\_savings\_egp\_per\_kwh} is
$S_{\text{ann}} / \sum_m G_m$ --- the \emph{effective} per-kWh
savings rate the household actually realises, which under the
EgyptERA schedule can be 2--3$\times$ the household's flat-tariff
\emph{average} rate. This gap is the numerical backbone of
Contribution~B: a flat-tariff PV calculator systematically
under-counts Egyptian payback by exactly this factor.

\subsection{PV-size optimisation under the tier schedule}

A flat-tariff calculator says ``more generation is always better''.
Under the marginal block tariff the optimum is finite: once
consumption is pulled down into the cheap tiers, additional
generation has near-zero value (because the export credit is zero).
The optimiser sweeps a candidate grid of system sizes
$\{P_k\}_{k=1}^{N_{\text{cand}}}$ (linear sweep between
user-supplied bounds), evaluates each candidate with the deterministic
financial kernel (\S\ref{sec:finance}), and returns the size that
maximises NPV:
\begin{equation}
P^{\star} = \arg\max_{P_k}\, \mathrm{NPV}(P_k).
\label{eq:optimum}
\end{equation}

The full set of candidates is returned alongside $P^{\star}$ so the
methodology section's NPV-vs-size curve can be plotted directly from
the API response.

% =====================================================================
\section{Financial Modelling}
\label{sec:finance}
% =====================================================================

\subsection{Year-by-year cash-flow chain}

Year 0 is the moment of installation (capex paid). Year $t \ge 1$ is
the $t$-th full year of operation. For an analysis horizon of $T$
years, year-1 generation $E_1$, year-1 effective tariff $p_1$, system
size $P_{\text{sys}}$, installed cost $C_{\text{kW}}$ (EGP/kW),
discount rate $r$, tariff escalation rate $i$, module degradation
rate $d$, and O\&M as a fraction $f$ of capex:
\begin{equation}
\begin{aligned}
\text{Capex}      &= P_{\text{sys}} \cdot C_{\text{kW}}                   \\
E_t               &= E_1 \cdot (1 - d)^{t-1}                              \\
p_t               &= p_1 \cdot (1 + i)^{t-1}                              \\
S_t               &= E_t \cdot p_t                                        \\
\text{O\&M}_t     &= f \cdot \text{Capex}                                 \\
\text{Net}_t      &= S_t - \text{O\&M}_t                                  \\
\mathrm{NPV}      &= -\,\text{Capex} + \sum_{t=1}^{T}
                     \frac{\text{Net}_t}{(1+r)^t}                         \\
\mathrm{LCOE}     &= \frac{\text{Capex}
                          + \sum_{t=1}^{T} \text{O\&M}_t / (1+r)^t}
                         {\sum_{t=1}^{T} E_t / (1+r)^t}
\end{aligned}
\label{eq:cashflow}
\end{equation}

The simple payback is the textbook
$\text{Capex} / (S_1 - \text{O\&M}_1)$, capped at $T$ if the project
fails to recover within the horizon. The \emph{discounted} payback is
the year at which the discounted cumulative cash flow first turns
non-negative, with linear interpolation between bracketing years:
\begin{equation}
t^{\star} = \min\bigl\{\, t : \mathrm{CCF}_t \ge 0 \,\bigr\},
\quad
\hat{t} = (t^{\star} - 1)
       + \frac{|\mathrm{CCF}_{t^{\star} - 1}|}
              {\mathrm{CCF}_{t^{\star}} - \mathrm{CCF}_{t^{\star} - 1}},
\quad
\mathrm{CCF}_t = -\,\text{Capex}
              + \sum_{s=1}^{t} \frac{\text{Net}_s}{(1+r)^s}.
\label{eq:payback}
\end{equation}

\paragraph{Why both simple and discounted payback?}
Simple payback is widely criticised in the energy-finance literature
for ignoring the time value of money and the post-payback period,
yet it is the single number every PV brochure quotes. A defensible
thesis must report it (so the result is comparable to the literature
and to consumer-facing tools) \textbf{and} report
NPV/LCOE/discounted payback alongside, so the reader can see the gap
that the na\"{i}ve metric hides. The kernel returns both.

\subsection{Egypt-specific economic constants}

The default values for $r, i, d, f, C_{\text{kW}}, T$ are listed in
\S\ref{sec:constants}. The 25-year horizon $T$ matches the standard
PV module performance warranty and is the analysis period adopted in
most Egyptian PV pre-feasibility studies \citep{mahmoud2023}.

% =====================================================================
\section{Uncertainty Quantification --- Monte Carlo (Contribution C)}
\label{sec:montecarlo}
% =====================================================================

The deterministic financial kernel returns a \emph{point} estimate of
payback, NPV, and LCOE. The third thesis contribution treats the
seven parameters with the largest real-world uncertainty as
\emph{probability distributions} and re-evaluates the same kernel
$N_{\text{sim}}$ times to produce \emph{distributions} of the
headline metrics. The default $N_{\text{sim}} = 1\,000$ delivers
$\pm 0.05$-year confidence on the median payback by Hoeffding's
inequality at the observed sample standard deviation.

\subsection{Parametric distributions and priors}

Each uncertain input gets a prior anchored in the published
literature. Two families are supported (normal and triangular),
because they cover the academically relevant cases without forcing
the API surface to grow as new parameters are added.

\begin{table}[!htbp]
\centering
\footnotesize
\setlength{\tabcolsep}{4pt}
\renewcommand{\arraystretch}{1.2}
\begin{tabularx}{\linewidth}{@{}p{3.2cm} l X@{}}
\toprule
\textbf{Parameter} & \textbf{Distribution} & \textbf{Source / justification} \\
\midrule
Module degradation $d$
  & $\mathrm{Tri}(0.002, 0.005, 0.010)$
  & NREL bounds for mono-Si \citep{jordan2013} \\
Tariff inflation $i$
  & $\mathrm{N}(0.08, 0.03)$, clip $\ge 0$
  & EgyptERA decade trend \citep{egyptera_tariff} \\
O\&M fraction $f$
  & $\mathrm{Tri}(0.005, 0.010, 0.020)$
  & IRENA residential rooftop \citep{irena2023} \\
Installed cost $C_{\text{kW}}$
  & $\mathrm{Tri}(30\text{k}, 35\text{k}, 45\text{k})$
  & Egyptian market 2024, EGP/kW \\
Annual yield factor
  & $\mathrm{N}(1.0, 0.05)$, clip $[0.5, 1.5]$
  & Egyptian PV field studies (TMY residual) \\
Inverter replacement year
  & $\mathrm{Tri}(10, 12, 15)$
  & IEA-PVPS Task 13 \citep{ieapvps2021} \\
Inverter cost fraction
  & $\mathrm{Tri}(0.07, 0.10, 0.15)$
  & Egyptian installer quotes \\
\bottomrule
\end{tabularx}
\caption{Prior distributions for the seven Monte Carlo inputs.
$\mathrm{N}(\mu,\sigma)$ denotes Normal, $\mathrm{Tri}(a,m,b)$
denotes Triangular with mode $m$.}
\label{tab:priors}
\end{table}

\paragraph{Why parametric and not bootstrap?}
A homeowner deciding on PV at $t = 0$ has no historical sample of
\emph{their own} future panel performance --- only published priors
on degradation, on tariff policy, on weather. The right uncertainty
model is therefore parametric, with priors anchored in the
literature. The kernel exposes every distribution as a request-level
override so a sensitivity-aware caller can collapse any distribution
to a deterministic constant or substitute a site-specific prior.

\paragraph{Why per-year, per-simulation yield noise?}
Annual irradiance in Egypt varies ${\sim}\pm 5\,\%$ around the TMY
mean (Egyptian PV field studies), and that variability \emph{does
not} average out over the analysis horizon for the \emph{payback}
metric --- early bad years push payback later in nonlinear ways. The
kernel therefore samples a fresh yield factor for each
$(\text{simulation}, \text{year})$ pair, giving a
$(N_{\text{sim}}, T)$ matrix.

\paragraph{Why one inverter replacement event, not zero or two?}
Modern string inverters in Egypt's climate carry 10--12 year
warranties and see a typical 12--15 year service life
\citep{ieapvps2021}. A 25-year analysis horizon therefore captures
one replacement; modelling two would over-attribute uncertainty to
the inverter alone and is left for the sensitivity tornado
(\S\ref{sec:tornado}).

\subsection{Sampling and aggregation}

Sampling uses a NumPy \texttt{Generator} with an optional
caller-supplied \texttt{random\_seed} so any reported figure is
byte-reproducible. Each distribution is drawn with the requested
shape (scalar per simulation, except for
\texttt{annual\_yield\_factor} which is per
$(N_{\text{sim}}, T)$); clipping is applied as a last step. The
cash-flow simulation is fully vectorised in NumPy along the
simulation axis: a $1\,000\times25$ ensemble runs in under
$100~\mathrm{ms}$ on a laptop, which matters because the dashboard
re-runs the engine on every user input change.

The reported aggregates per metric are
$\{\mathrm{mean}, \mathrm{std}, p_{05}, p_{10}, p_{25}, p_{50},
p_{75}, p_{90}, p_{95}, \min, \max\}$. Reporting both
$\mathrm{mean} \pm \mathrm{std}$ and the 5--95\,\% percentile band
makes the asymmetry of the simulated distribution visible: payback in
particular is heavily right-skewed (a fat tail of ``never recovers''
runs) and a Gaussian-style $\mathrm{mean} \pm \mathrm{std}$ would
mis-represent it on its own.

\subsection{Cumulative-cash-flow trajectory bands (Day-16 fan chart)}

Beyond per-metric percentiles, the engine also returns the
\emph{year-by-year} percentile bands of the discounted cumulative
cash flow. For each year $t \in \{0, 1, \dots, T\}$ and percentile
$q \in \{5, 25, 50, 75, 95\}$,
\begin{equation}
\mathrm{CCF}^{(q)}(t) =
   \mathrm{percentile}_q
   \bigl\{ \mathrm{CCF}^{(k)}(t) \,:\, k = 1\dots N_{\text{sim}}\bigr\}.
\label{eq:fanchart}
\end{equation}
The bands are \emph{envelope} percentiles, not the trajectory of any
single simulation. Reporting bands rather than a single mean curve is
the whole point of the contribution: a homeowner sees not just ``the
median crosses zero in year 7'' but also ``the worst-case 5\,\% of
futures still owe money in year 12'' and ``the best 5\,\% of futures
double their money by year 15''.

% =====================================================================
\section{Sensitivity Tornado (Day 18)}
\label{sec:tornado}
% =====================================================================

The tornado is the \emph{attribution} counterpart to the Monte Carlo
\emph{total-uncertainty} view. For each input parameter $\theta_p$
($p = 1\dots P$), one keeps every other parameter at its baseline
and re-runs the deterministic financial kernel at
$\theta_p \in \{\theta_p^{\text{low}}, \theta_p^{\text{high}}\}$,
recording the resulting metric values
$M_p^{\text{low}}, M_p^{\text{high}}$. The tornado bar for parameter
$p$ has length $|M_p^{\text{high}} - M_p^{\text{low}}|$; rows are
sorted by that absolute swing descending.

The seven parameters swung --- \texttt{annual\_kwh},
\texttt{tariff\_egp\_per\_kwh}, \texttt{cost\_egp\_per\_kw},
\texttt{discount\_rate}, \texttt{tariff\_inflation\_rate},
\texttt{annual\_degradation\_rate}, \texttt{om\_cost\_fraction}
--- exactly mirror the Monte Carlo stochastic inputs
(\S\ref{sec:montecarlo}) plus the household-level levers (tariff,
generation) that turn project-level uncertainty into household-level
decision support. The swing endpoints default to literature-anchored
low/high values for each parameter; callers can supply explicit
ranges per parameter via the API surface.

\paragraph{Why OAT and not Sobol or variance decomposition?}
Sobol indices give a more complete picture of joint sensitivity but
are not directly interpretable as ``this parameter changes my NPV by
$\pm X~\mathrm{EGP}$'' --- they are \emph{fractional contributions to
output variance} and require a non-trivial statistical literacy to
read. The OAT tornado is the standard reporting format in the
rooftop-PV pre-feasibility literature
\citep{nrel-sam, ieapvps-task7} and the format the bachelor-thesis
dashboard's homeowner audience can read in a single pass. The Day-9
Monte Carlo engine, which \emph{does} model joint uncertainty, is
the complementary figure --- together they cover both sensitivity
questions a methodology section is expected to address.

% =====================================================================
\section{\texorpdfstring{$\COtwo$}{CO2} Avoidance (Day 18)}
\label{sec:co2}
% =====================================================================

For year-1 generation $E_1$, analysis horizon $T$, degradation rate
$d$, and grid emission factor $\varepsilon_{\text{grid}}$
(EEHC 2022/2023, $0.46~\mathrm{kg\,\COtwo/kWh}$):
\begin{equation}
(\COtwo)_t = E_1 \cdot (1 - d)^{t-1} \cdot \varepsilon_{\text{grid}},
\qquad
(\COtwo)_{\text{lifetime}} = \sum_{t=1}^{T} (\COtwo)_t.
\label{eq:co2}
\end{equation}

\paragraph{Why a marginal grid-average emission factor and not a
true marginal-dispatch factor?}
The Egyptian grid's marginal-dispatch emission factor (the kg
$\COtwo$ avoided by displacing the \emph{next} kWh the grid would
have generated) varies by hour, season, and merit-order conditions.
Public, peer-reviewed Egyptian time-resolved marginal data is not
available, so the standard practice in Egyptian PV pre-feasibility
literature \citep{mahmoud2023} is to use the EEHC published
\emph{grid-average} annual emission factor. This biases the result
conservatively whenever PV displaces high-merit gas peakers --- the
marginal factor is typically 10--20\,\% higher than the grid-average
in gas-dominated systems. The kernel exposes the factor as an
override so a methodology-aware user can substitute a marginal
number when one becomes available.

\paragraph{Why no embodied-carbon subtraction?}
A complete LCA would net off the embodied carbon of the modules,
inverter, and balance-of-system (typically
30--50\,$\mathrm{g\,\COtwo/kWh}$ amortised over a 25-year horizon,
IEA-PVPS Task 12 \citep{ieapvps2020}). Including a half-modelled LCA
in the headline number would over-claim precision the dataset does
not support; we flag this in \texttt{limitations.tex} as future work.

\paragraph{Homeowner-friendly equivalences.}
The lifetime kg figure is also reported as
\begin{itemize}[leftmargin=*]
\item equivalent passenger-car kilometres at $0.12~\mathrm{kg/km}$
(EEA fleet-average tail-pipe);
\item equivalent litres of petrol burnt at $2.31~\mathrm{kg/L}$
(EPA Greenhouse Gas Equivalencies Calculator);
\item equivalent urban trees grown over the analysis horizon at
${\sim}21~\mathrm{kg\,\COtwo/tree/yr}$ (US EPA midpoint).
\end{itemize}
The two anchors have very different orders of magnitude so no single
equivalence dominates the reader's intuition.

% =====================================================================
\section{Reproducibility}
% =====================================================================

All deterministic kernels (sizing, energy, financial, tariff,
$\COtwo$, sensitivity) are pure functions of their request payloads
--- no hidden state, no environment variables consumed inside the
kernel. The Monte Carlo engine accepts a \texttt{random\_seed} that
threads into the underlying NumPy \texttt{Generator}; the test suite
relies on this for byte-identical re-runs and the API surface
exposes it so any reported figure in this thesis can be reproduced
with a single documented \texttt{curl} invocation.

External-API services (PVGIS, Overpass, Google Maps Static) are
hidden behind small adapter modules with explicit retry, timeout,
and error-translation behaviour. The full unit-test suite (358
tests as of Day 19) mocks every external call so a clean checkout
runs end-to-end without network access; the validation chapter
(Day 20) re-runs the same suite against live PVGIS for ten
representative Egyptian sites and reports the residuals.

The runtime environment is pinned to Python 3.12 with the
\texttt{requirements.txt} shipping pinned versions for
\texttt{pvlib}, \texttt{pandas}, \texttt{numpy}, \texttt{scipy},
\texttt{pillow}, \texttt{shapely}, \texttt{httpx}, \texttt{fastapi},
and \texttt{pydantic}. The full thesis result set can be re-derived
from the repository at any commit by running

\begin{lstlisting}[language=bash]
cd backend
python3 -m venv .venv && .venv/bin/pip install -r requirements.txt
.venv/bin/uvicorn app.main:app --reload
# ... then the documented curl invocations in each chapter
\end{lstlisting}

% =====================================================================
\section{Summary}
% =====================================================================

The four academic claims of the thesis decompose along the chapter
layout above:

\begin{table}[!htbp]
\centering
\footnotesize
\begin{tabularx}{\linewidth}{@{}X l@{}}
\toprule
\textbf{Claim} & \textbf{Section(s)} \\
\midrule
Dual energy model is the methodological backbone
  & \S\ref{sec:pvlib}, \S\ref{sec:manual} \\
AI-assisted roof detection eliminates the most error-prone user input
  & \S\ref{sec:osm}, \S\ref{sec:cv} \\
Egypt's tiered tariff is what a flat-tariff calculator gets wrong
  & \S\ref{sec:tariff-blocks}--\ref{sec:tariff} \\
Monte Carlo + tornado make uncertainty an artefact, not a postscript
  & \S\ref{sec:montecarlo}, \S\ref{sec:tornado} \\
\bottomrule
\end{tabularx}
\caption{Mapping of academic claims to chapter sections.}
\label{tab:claims}
\end{table}

Every constant in \S\ref{sec:constants} is exposed as a
Pydantic-validated override so a methodology-aware caller can
substitute a site-specific value without touching the source. Every
kernel is pure and deterministic (except the Monte Carlo engine,
which accepts a seed). Every figure in this chapter resolves either
to one of the seven backend service modules or to one of the five
frontend chart components --- a one-to-one mapping that makes the
methodology auditable end-to-end.
